\section{Challenges in Predictive Modeling of Critical Care: Addressing Missing Features and Limited Labels}
Recent research has increasingly concentrated on utilizing time-series data of laboratory results and vital signs to predict patient outcomes in critical care settings. This chapter explores two specific challenges associated with these predictive models.
First, \emph{feature missingness is common in patient records either due to varied frequency of measurements of various labs and vitals, or if certain physiological measurements are deemed irrelevant for diagnosis.}
Second, for many diagnostic tasks, obtaining outcome labels to supervise training is prohibitively expensive, leading to \emph{small labeled sets}.
Both challenges require methods that inherently address uncertainty.
We argue that simple probabilistic models -- hidden Markov models (HMMs)~\citep{rabinerIntroductionHiddenMarkov1986,ghahramaniIntroductionHiddenMarkov2001b} with a few dozen states -- have been under-explored yet provide a natural solution. Our experiments show that HMMs with 5-20 states, requiring minimal computational resources, are capable of elegantly capturing temporal trends across various physiological channels, resulting in prediction performance competitive with large sequential neural nets.

To address the problem of \emph{feature missingness}, most efforts to date targeted at critical care time series consider two kinds of approaches: (1) a preprocessing step that imputes missing values via a simple heuristic like population mean filling or forward filling~\citep{purushothamBenchmarkingDeepLearning2018,harutyunyanMultitaskLearningBenchmarking2019a, harvey2023comparative}, followed by a standard time series classifier, or (2) a deep learning approach that performs iterative refinement of missing values following an initial guess~\citep{cheRecurrentNeuralNetworks2018,yoonGAINMissingData2018,caoBRITSBidirectionalRecurrent2018}.
Unlike these approaches, HMMs are probabilistic models that elegantly handle missingness via exact marginalization, avoiding any imputation on the way to prediction.
Recent work has also explored other \emph{deep probabilistic} approaches to missingness that hybridize neural nets with Gaussian processes~\citep{futomaLearningDetectSepsis2017},
latent or stochastic differential equations~\citep{rubanovaLatentODEsIrregularlySampled2019,debrouwerGRUODEBayesContinuousModeling2019},
Hawkes processes~\citep{meiNeuralHawkesProcess2017},
 or encoder-decoder models~\citep{liLearningIrregularlySampledTime2020}.
Unlike these approaches, HMMs do not need tens of thousands of parameters or specialized GPU hardware to make good predictions. 

To address the \emph{small labeled set} problem, one promising approach is to use semi-supervised learning (\emph{SSL})~\citep{vanengelenSurveySemisupervisedLearning2020}, which trains models on the union of a small labeled set and a large unlabeled set of only observed time-series of features, without any known labels.
While applications of SSL to critical care time series are limited \citep{ahuja2022semi, poulain2022few, lam2021semisupervised}, there are many tasks for which it would be applicable, such as automating preliminary screening of diseases such as sepsis. One line of work has produced effective SSL for end-to-end deep classifiers of \emph{images}~\citep{miyato2018virtual,berthelotMixMatchHolisticApproach2019,sohnFixMatchSimplifyingSemisupervised2020}.
Recently, \citet{goschenhofer2021deep} suggest these methods can be effective for multivariate time-series classification as well, assuming no features are missing. 
A few very recent methods \citep{shin2023context,xing2022selfmatch} have developed SSL approaches specialized to time series classification, but neither can handle any missingness in the time series features.
Unlike such deep SSL classifiers, we show how HMMs as generative models can be trained on both small labeled sets and large unlabeled sets with arbitrary feature missingness.

\textbf{Contribution 1: HMMs as a feature extractor for time series with missingness.}
We show how HMMs can produce a learned fixed-dimensional representation vector from any regularly-sampled time series, regardless of missingness or length.
We take advantage of the probabilistic nature of HMMs, estimating the representation conditioning only on the observed features.
This effectively \emph{marginalizes} (integrates away) unobserved feature dimensions.
In practical terms, this means we only evaluate the each timestep's likelihood at observed dimensions.
There is no direct imputation or approximation of any missing feature values, eliminating a source of bias or noise.

\textbf{Contribution 2: Effective supervision and semi-supervision for HMMs.} 
Our experiments find that common unsupervised or supervised approaches to training latent variable models like HMMs for prediction tasks yield sub-par predictions.
To make HMMs succeed at clinical time-series prediction, especially with limited labels, our key methodological insight is \emph{prediction constrained} (PC) training~\citep{hopePredictionConstrainedHiddenMarkov2021,hughes_prediction-constrained_2017}.
This optimization objective balances the HMM's useful generative properties with the essential need for its learned representations to enable effective prediction in the clinical task. 
We show that PC training outperforms two common alternatives. First, the two-stage pipeline of unsupervised training of an HMM followed by supervised fine-tuning of a classifier on the frozen HMM's features (to maximize the likelihood of labels given features).
Second, supervised training of an HMM to maximize the joint likelihood of features and labels.

To highlight the versatility of the prediction constrained HMM (PC-HMM) for clinical prediction tasks, we perform experiments pursuing a binary classification task (in-ICU mortality prediction) and an ordinal classification task (categorizing length of stay in the ICU) on two large health records datasets.
We show that PC training of HMMs of modest size can match or outperform much larger deep neural networks designed specifically to match one of our stated challenges (feature missingness, limited labels).
We compare to 
GRU-D~\citep{cheRecurrentNeuralNetworks2018} and BRITS \citep{caoBRITSBidirectionalRecurrent2018} as strong baselines for missing features.
In the limited labels regime, we compare time-series adaptations~\citep{goschenhofer2021deep} of two strong semi-supervised methods for deep classifiers: MixMatch~\citep{berthelotMixMatchHolisticApproach2019} and FixMatch \citep{sohnFixMatchSimplifyingSemisupervised2020}.
We further show how compact HMMs can be \emph{interpreted} to yield useful clinical insights red both at an individual patient level and over the entire population.
Our PC-HMM results suggest that simple, easy-to-train probabilistic approaches deserve consideration in clinical applications, especially when computational resources are limited.
From the environmental costs of training deep models~\citep{strubellEnergyPolicyConsiderations2019} to avoiding expensive GPUs in bespoke hospital settings,
there are many reasons to try the PC-HMM in clinical time-series applications with missing features or labels. 

\subsection{Benefits Of Our Generative Modeling Approach}
Generative models learn the joint probability distribution $P(X,Y)$ of the input data $X$ and the output labels $Y$. These models attempt to model how the data is generated by capturing the underlying distribution of the data itself. Discriminative models, on the other hand, learn the conditional probability distribution $P(Y|X)$. They focus on the boundary that separates different classes and are directly optimized for classification tasks.
In this work, we choose the generative modeling approach for 2 main reasons:
\begin{itemize}
    \item Generative models offer a probabilistic framework that effectively manages missing features through exact marginalization.

    \item Generative models are ideal for leveraging both labeled and unlabeled data, maximizing the use of the large unlabeled set and improving predictive performance even when labeled data is scarce.
\end{itemize}
Later in our experiments, we also show that our generative model is capable of achieving better predictive performance than several discriminative model alternatives. 

\subsection{Is Our Model Implicit or Explicit?}
Our proposed Prediction Constrained Hidden Markov Model (PC-HMM) is an explicit generative model. 

\textbf{Implicit Generative Models \citep{xie2016theory, ngiam2011learning} :} Implicit generative models do not specify a direct form for the probability distribution of the data. Instead, they generate samples from the distribution through a process that is not explicitly defined in terms of a probability density function. The training process focuses on improving the generator to produce increasingly realistic samples, but it does not result in an explicit probability distribution that can be easily sampled or evaluated. One example of such models is a Generative Adversarial Network (GAN). 


\textbf{Explicit Generative Models \citep{goodfellow2014generative, arjovsky2017wasserstein}:} Explicit generative models define a specific form for the probability distribution of the data, in terms of a joint probability distribution \(P(X, Z)\), where \(X\) represents the observed data and \(Z\) represents the hidden variables or states. These models use parameters to describe the distribution and provide a way to calculate the probability of observed data given the parameters.

\textbf{HMM as an Explicit Model} In the case of HMMs, they explicitly model the joint distribution of the observed sequence and the sequence of hidden states. An HMM consists of:

\begin{itemize}
    \item Transition Probabilities: \(P(z_t | z_{t-1})\) - The probability of transitioning from one hidden state \(z_{t-1}\) to another hidden state \(z_t\).

    \item Emission Probabilities: \(P(x_t | z_t)\) - The probability of observing \(x_t\) given the hidden state \(z_t\).

    \item Initial State Distribution: \(P(z_1)\) - The probability distribution over initial states.
\end{itemize}
These components of the joint distribution of observed data and hidden states allow for explicit computation of likelihoods and probabilistic inference, making our proposed model an explicit generative model. Our choice of explicit generative models aligns with our overall goal due to the following reasons : 
\begin{itemize}
    \item \textbf{Handling missing data :} Explicit generative models are particularly adept at managing missing data by specifying a model for the underlying data distribution. Implicit models, while they can handle missing data, often do so less effectively as they might not model the data distribution as comprehensively.

    \item \textbf{Interpretability :} The parameters in explicit models often have meaningful interpretations, such as emissions probabilities in Hidden Markov Models (HMMs) representing the likelihood of generating features given the latent state. This transparency helps clinicians trust and understand the model’s decisions. Implicit models might lack this transparency as they do not explicitly define the data distribution, making it harder to understand the underlying mechanisms.

    \item \textbf{Likelihood estimation :} Explicit generative models allow for direct likelihood estimation, which can be useful for model evaluation and comparison. Likelihood provides a measure of how well the model fits the observed data, facilitating model selection and tuning. Implicit models typically do not provide likelihood estimates, making it more challenging to quantitatively assess model performance.
\end{itemize}